\section{Arquitecturas alternativas, más allá de 1T e historia de la evolución de los modelos}

La segunda mitad de la entrevista repasa la evolución de los modelos en los últimos tres años: ChatGPT abrió la interacción conversacional de contexto corto; en 2023 los equipos de código abierto se pusieron al día en preentrenamiento; en la etapa de reasoning, el código y las matemáticas se volvieron señales verificables; en la etapa de Agent, el código y la retroalimentación del entorno se convirtieron en la vía de entrada para entrenar tareas de largo alcance. Luo Fuli (罗福莉) considera que la brecha generacional de los equipos chinos en pre-train ya es muy pequeña y que la próxima competencia está en el post-train de Agent y en el RL scaling.

\begin{figure}[H]
\centering
\includegraphics[width=0.96\textwidth]{model-evolution-timeline.png}
\caption{Evolución de los modelos en los últimos tres años: de ChatGPT al postentrenamiento de Agent.}
\end{figure}

\begin{knowledgebox}{Lectura de la figura: cómo cada etapa hereda de la anterior}
La línea de tiempo no es una relación de sustitución. ChatGPT permitió a los usuarios percibir la inteligencia del preentrenamiento; la recuperación del código abierto completó los datos y la arquitectura; reasoning hizo de las tareas verificables una señal de entrenamiento; los prototipos de Agent incorporaron el código y el entorno a las tareas de largo alcance; y 2026 llevó el post-train scaling a la línea principal.
\end{knowledgebox}

\subsection{El sentido de las arquitecturas alternativas}

Luo Fuli menciona decisiones «a contracorriente» como DeepSeek-V2, MoE y las modificaciones de attention. Una arquitectura alternativa no va contra la corriente por el simple hecho de hacerlo, sino que busca investigación escalable en entornos con recursos limitados o sensibles al costo. Si una investigación es solo un paper trick, no necesariamente se convierte en un modelo de nivel industrial; pero si logra escalar y termina por alcanzar un modelo de alto nivel, es frontier research valiosa.

\noindent\begin{tabularx}{\textwidth}{p{0.22\textwidth}p{0.34\textwidth}X}
\toprule
Ruta & Problema que resuelve & Criterio de industrialización \\
\midrule
MoE & Reducir el costo de cómputo por token mediante activación dispersa & Requiere que el enrutamiento, el balanceo de carga, la estabilidad del entrenamiento y la eficiencia del despliegue se cumplan a la vez. \\
Modificación de attention & Reducir el costo de contexto largo y de inferencia & Hay que ver si mantiene la calidad en tareas reales, no solo en un benchmark local. \\
Modelo pequeño + marco & Usar un modelo barato junto con un Agent scaffold & Adecuado para etapas de alta frecuencia, bajo costo y cuyas carencias puede suplir el marco. \\
Base de 1T & Proporcionar la base de capacidades de Agent de frontera & Es el boleto de entrada, no la respuesta completa. \\
\bottomrule
\end{tabularx}

\subsection{La IA no tiene crisis de supervivencia}

Luo Fuli compara la evolución humana con la evolución de los modelos: los seres humanos evolucionaron en un entorno natural y bajo presión de supervivencia; los modelos no tienen una crisis de supervivencia propia, pero cuentan con mucho cómputo, con el conocimiento humano y con la ayuda de las personas. Por ello, la trayectoria evolutiva de la IA puede ser más libre, más dispersa, más rápida y también más distinta de la inteligencia biológica. Este pasaje nos recuerda que no debemos aplicar de forma forzada la analogía de la evolución humana a los modelos.

\begin{knowledgebox}{Los límites de la analogía}
«El lenguaje apareció muy tarde en la inteligencia humana» puede ayudarnos a entender la estructura de triángulo invertido de la IA, pero el entorno de los modelos es completamente distinto del humano. El siguiente paso de la IA no necesariamente reproducirá la trayectoria del cuerpo o de los sentidos humanos; es posible que primero siga expandiéndose en coding, en productividad y en entornos digitales.
\end{knowledgebox}

\subsection{Negar cada día al yo de ayer}

Luo Fuli dice que en el último medio año casi todos los días ha negado a su yo del día anterior. Esta expresión describe bien la etapa de Agent de 2026: las herramientas, los marcos, los modelos, los costos y las formas de organización cambian con rapidez. Para un equipo de investigación, la estabilidad no consiste en aferrarse a juicios antiguos, sino en construir mecanismos que permitan actualizarlos con rapidez.

\begin{importantbox}{De dónde viene la velocidad de investigación}
La velocidad de investigación no proviene solo de personas inteligentes, sino también del entorno: cómputo suficiente, infraestructura suficientemente buena, retroalimentación suficientemente rápida, una organización suficientemente abierta y valores suficientemente claros. Que el entorno importe más que la experiencia es el juicio central con el que cierra este episodio.
\end{importantbox}

\subsection{Resumen de la sección}

En los últimos tres años, la evolución de los modelos pasó de la interacción conversacional, la recuperación del código abierto y la verificación mediante reasoning al postentrenamiento de Agent. El liderazgo en la siguiente etapa dependerá de arquitecturas que puedan escalar, de los sistemas de postentrenamiento, de la asignación de recursos y de la capacidad de negarse a sí mismo con rapidez.
